
\subsubsection{3.1 Dataset}

A dataset of 120 Papers with Code benchmark repositories was collected using the GitHub REST API v3. Repositories were selected with minimum 32 stars and excluded forks to ensure original projects. The curated list included repositories across computer vision, natural language processing, reinforcement learning, general machine learning, MLOps, and data processing domains.

Six core GitHub metadata features were extracted per repository:
\begin{itemize}
\item stars\_log: $log_1$₊(stargazers\_count)
\item forks\_log: $log_1$₊(forks\_count)
\item contributors\_log: $log_1$₊(contributor count)
\item total\_commits\_log: $log_1$₊(commit count)
\item open\_issues\_log: $log_1$₊(open\_issues\_count)
\item days\_since\_last\_commit: current date minus last commit timestamp
\end{itemize}

$Log_1$₊ transformation (log(1+x)) was applied to address long-tail distributions in GitHub metadata.

Binary labels were assigned using a 180-day threshold: maintained (days\_since\_last\_commit < 180), abandoned (days\_since\_last\_commit $\geq$ 180). This threshold follows He et al. (2024).

Dataset statistics: 120 total repositories, 99 maintained (82.5\%), 21 abandoned (17.5\%). Stratified 80/20 train/test split yielded 96 training samples (79 maintained, 17 abandoned) and 24 test samples (20 maintained, 4 abandoned). Random seed 42 was used for reproducibility.

Example repositories included huggingface/transformers (120k stars, 1 day since last commit), pytorch/pytorch (76k stars, 1 day), apache/mxnet (20.6k stars, 450 days), and deepmind/acme (3.4k stars, 180 days).

Dataset size rationale: the target was 2000 repositories, but GitHub API rate limits (60 unauthenticated requests/hour) constrained collection to 120 repositories. All repositories were real, verified, and accessible on GitHub.

\subsubsection{3.2 Models}

Two models were trained in sequence: logistic regression as the simplicity test, then gradient boosting for complexity comparison.

Logistic regression used sklearn.linear\_model.LogisticRegression with parameters: max\_iter=1000, class\_weight='balanced' (to handle 82.5\% vs 17.5\% imbalance), solver='lbfgs', random\_state=42, C=1.0 (L2 regularization). StandardScaler normalization (zero mean, unit variance) was fitted on the training set and applied to training and test sets. Training time was approximately 30 seconds on a single CPU. Convergence occurred in 16 iterations.

Gradient boosting used sklearn.ensemble.GradientBoostingClassifier with parameters: n\_estimators=50, max\_depth=6, learning\_rate=0.1, random\_state=42. Class imbalance was handled via sample weighting. Training time was approximately 10 minutes on a multi-core CPU.

\subsubsection{3.3 Evaluation}

Performance metrics included accuracy, precision, recall, F1 score, and ROC-AUC. Test set performance (24 samples) was the primary evaluation.

Hypothesis validation gates were defined as follows. H-E1 (existence): logistic regression achieves accuracy $\geq 75$\% AND F1 $\geq 0.73$. H-M1 (mechanism): (a) logistic regression coefficient signs match causal predictions (days\_since\_last negative, activity features positive), (b) logistic regression -- gradient boosting performance gap $\leq 5$\%, (c) feature importance overlap $\geq 2/3$ (at least 2 of top-3 features agree between models).

Five visualizations were generated: confusion matrix, ROC curve, coefficient bar chart, feature importance comparison, and PCA-projected decision boundary.
